\section{Capítulo~\ref{sec:synthesis}. La máquina completa}

El capítulo de síntesis no presenta experimentos nuevos: cada fila se apoya en el capítulo donde se analiza la proposición y en las mismas fuentes; el bus \nt{GLUT} y el control del flujo mediante \nt{GABA} están firmemente establecidos, mientras que los papeles de los canales de difusión y su trabajo conjunto son sobre todo una hipótesis.

{\footnotesize
\setlength{\tabcolsep}{3pt}\renewcommand{\arraystretch}{1.15}
\begin{longtable}{>{\raggedright}p{0.5cm}>{\raggedright}p{4.0cm}>{\raggedright}p{5.6cm}>{\raggedright}p{2.3cm}>{\raggedright\arraybackslash}p{1.7cm}}
\toprule
N.º & Afirmación & Experimento y qué se midió & Fuente & Estado \\
\midrule\endfirsthead
\toprule
N.º & Afirmación & Experimento y qué se midió & Fuente & Estado \\
\midrule\endhead
1 & Para $8.6\cdot10^{10}$ neuronas hay solo cuatro clases de señales de control~\eqref{eq:machine} & Recuento de núcleos celulares en un homogeneizado de cerebro (fraccionador isotrópico): un hombre adulto tiene $86.1\pm8.1$ mil millones de neuronas & \cite{Azevedo2009} & medido \\
2 & Proposición~\ref{prop:architecture}, las dos primeras capas: los datos viajan por el bus de direcciones \nt{GLUT}, el signo de la conexión lo fija el emisor, y las neuronas \nt{GABA} guían y normalizan el flujo (capítulos~\ref{sec:neurontypes}, \ref{sec:gaba}) & Un neurotransmisor principal por neurona (revisión de mediciones); la inhibición prealimentada estrecha la ventana en la que la neurona piramidal suma sus entradas; la división por la actividad total se ha medido en muchas áreas & \cite{Strata1999,Pouille2001,Carandini2012} & medido \\
3 & Los elementos no son fiables: la sinapsis pierde la mayoría de las espigas (tabla~\ref{tab:computer}, capítulo~\ref{sec:code}) & Cortes de hipocampo, estimulación mínima: más de la mitad de las espigas no provoca respuesta en CA1; se descartan fallos del umbral y de la conducción por el axón; la causa es la liberación probabilística del neurotransmisor & \cite{AllenStevens1994} & medido \\
4 & El cerebro funciona con $\sim20$ vatios, en promedio cerca de una espiga por segundo por neurona (capítulo~\ref{sec:code}); los ingenieros aún no han construido una máquina así & Estimación~\eqref{eq:energy} a partir de los costos medidos: unas $10^9$ moléculas de ATP por espiga, incluido el trabajo de las sinapsis (corteza de roedores); una estimación detallada para la corteza da una frecuencia aún menor. Estado de la técnica, según una revisión & \cite{Attwell2001,Lennie2003,MehonicKenyon2022} & teoría \\
5 & El aprendizaje de la mayoría de las sinapsis es el producto de una traza local por un único número común $\delta$ (segunda ecuación de~\eqref{eq:machine}, capítulo~\ref{sec:dofa}) & Las neuronas \nt{DOPA} del mono responden al error de predicción de la recompensa; en cortes de estriado, la dopamina agranda la espina solo si llega $0.3$--$2$\,s después de la entrada glutamatérgica: la traza está medida & \cite{Schultz1997,Yagishita2014} & medido \\
6 & Un cable de error propio para cada salida existe solo en el circuito del aprendizaje motor (capítulo~\ref{sec:motorlearning}) & Cerebelo: la coincidencia de la señal de la fibra trepadora con una entrada debilita esa entrada; en el mono, las espigas complejas de las células de Purkinje llevan la dirección del error del movimiento y provocan la corrección & \cite{Ito2001,Herzfeld2018} & medido \\
7 & Cada canal de difusión es un mazo de unas pocas líneas (proposición~\ref{prop:bundle}) & Para \nt{DOPA}: en una misma tarea, distintas neuronas llevan la recompensa, el movimiento y los rasgos del estímulo; el error rápido y el nivel lento de dopamina son distintos. Para \nt{SERO}, \nt{NORA} y \nt{ACET} este desglose está menos estudiado & \cite{Engelhard2019,Hamid2016,Mohebi2019} & medido \\
8 & \nt{SERO}: regulador de nivel y, según la hipótesis, horizonte $\tau_\gamma$ (capítulo~\ref{sec:serocomputer}) & Autoinhibición: los agonistas de sus propios receptores 5-HT$_{1A}$ inhiben las neuronas serotoninérgicas del rafe (medido). El horizonte se propuso en la teoría; el experimento más fuerte: la activación optogenética de estas neuronas en el ratón alarga la espera de una recompensa diferida & \cite{Sprouse1987,Doya2002,Miyazaki2014} & hipótesis \\
9 & \nt{NORA}: la ganancia $g$ elige entre buscar y decidir (capítulo~\ref{sec:nora}) & Aumento de la relación señal/ruido: en el mono despierto, la noradrenalina en la corteza auditiva suprime la actividad de fondo más que la respuesta a las voces de congéneres (medido); la ganancia multiplicativa es un modelo; el papel en la elección entre buscar y decidir, una teoría & \cite{Foote1975NE,ServanSchreiber1990,AstonJones2005} & hipótesis \\
10 & \nt{ACET}: alterna escritura y lectura (capítulo~\ref{sec:acet}) & Las tres acciones (debilitar las conexiones internas, reforzar la entrada, facilitar la plasticidad) están medidas; el cambio de modo lo apoya indirectamente un experimento con escopolamina en humanos & \cite{HasselmoBower1992,Gil1997,Atri2004} & hipótesis \\
11 & La fórmula~\eqref{eq:machine} describe la máquina completa & Es un resumen de los modelos de los capítulos~\ref{sec:glutcomputer}--\ref{sec:acet}; cada parte se apoya en su capítulo; como un todo no se ha puesto a prueba experimentalmente & deducción en el capítulo & hipótesis \\
12 & Las perillas trabajan de forma coordinada sin un control común: error, sorpresa, escritura, regreso (fig.~\ref{fig:timeline}) & No hay experimento: el esquema es cualitativo. Eslabones sueltos: la teoría del reparto de tareas entre \nt{NORA} y \nt{ACET} y la relación entre la pupila y la tasa de aprendizaje en humanos & \cite{YuDayan2005,Nassar2012} & hipótesis \\
13 & El aprendizaje con una única señal común es tanto más lento cuantas más neuronas influyen en el resultado (proposición~\ref{prop:broadcastcost}) & Cálculo de las curvas de aprendizaje de una red lineal: el aprendizaje por perturbación de las salidas es más lento que el descenso por gradiente tantas veces como salidas haya & \cite{Werfel2005} & teoría \\
14 & No se sabe cómo la corteza ajusta los circuitos de muchas capas; candidatos: ráfagas de espigas, cálculos en las ramas de la neurona, autoaprendizaje por predicción & Las ráfagas como portadoras del error: modelo; solo una red de $5$--$8$ capas logró reproducir la respuesta de un modelo detallado de neurona cortical: modelo; espigas de calcio en las ramas de neuronas corticales humanas que dan el OR exclusivo: medido en cortes; autoaprendizaje: revisión de pruebas & \cite{Payeur2021,Beniaguev2021,Gidon2020,Keller2018} & hipótesis \\
\bottomrule
\end{longtable}}
